\documentclass[aspectratio=169]{beamer}

\usetheme{Warsaw}
\usecolortheme{default}

\usepackage[utf8]{inputenc}
\usepackage{listings}
\usepackage{xcolor}
\usepackage{hyperref}
\usepackage{url}

% Configure code listings
\lstset{
    basicstyle=\ttfamily\footnotesize,
    breaklines=true,
    frame=single,
    showstringspaces=false,
    backgroundcolor=\color{gray!10},
    keywordstyle=\color{blue},
    commentstyle=\color{green!50!black},
    stringstyle=\color{red}
}


\definecolor{yamlkey}{rgb}{0.56,0.35,0.01}
\definecolor{yamlvalue}{rgb}{0.0,0.47,0.62}
\definecolor{yamlcomment}{rgb}{0.5,0.5,0.5}

\lstdefinelanguage{yaml}{
  keywords={true,false,null,y,n},
  keywordstyle=\color{blue},
  basicstyle=\ttfamily\small,
  sensitive=true,
  comment=[l]{\#},
  morecomment=[l]{\#},
  commentstyle=\color{yamlcomment}\ttfamily,
  morestring=[b]",
  stringstyle=\color{yamlvalue},
  moredelim=[s][\color{yamlkey}]{:}{\ },
}


\title{Omnibenchmark for solo benchmarking}
\author{Izaskun Mallona}
\date{\today}

\begin{document}

\frame{\titlepage}

\begin{frame}{Overview}
    \begin{itemize}
        \item What is Omnibenchmark?
        \item Architecture and design philosophy
        \item The clustering example
        \item How to run a benchmark
        \item Contributing modules
        \item Best practices
    \end{itemize}
\end{frame}

\begin{frame}{What is Omnibenchmark?}
    \textbf{Omnibenchmark} is a CLI to facilitate solo and collaborative benchmarking
    \textbf{Key features:}
    \begin{itemize}
        \item \textbf{Modular}: Distributed modules in separate repositories
        \item \textbf{Collaborative}: Community can contribute new modules
        \item \textbf{Transparent}: Topology, modules, metrics, and contribution are formally defined and reusable
        \item \textbf{Automated}: Workflow-manager friendly
        \item \textbf{Repeatable}: Includes several software backends (Conda, Apptainer, modules)
    \end{itemize}
    
    \vspace{0.3cm}
    Website: \url{https://omnibenchmark.org}
    
\end{frame}

\begin{frame}{Design philosophy}
    \textbf{Traditional benchmarks:}
    \begin{itemize}
        \item One-time publication
        \item Centralized, monolithic code
        \item Difficult to extend
        \item Results become outdated
    \end{itemize}
    
    \vspace{0.3cm}
    \textbf{Omnibenchmark approach:}
    \begin{itemize}
        \item Living benchmark: Continuously updated
        \item Distributed modules: Each in its own repository
        \item Easy extension: Add new modules without changing core
        \item Current results: Automatically rerun when updated
        \item Fair comparison: All methods use same infrastructure
    \end{itemize}
    
    \vspace{0.3cm}
    \textit{Think of it as a "continuous integration for benchmarking"}
\end{frame}

\begin{frame}{Benchmarks architecture}
  
   \textbf{Flexible topology, but often:}
    \begin{enumerate}
        \item \textbf{Data modules}: Generate or load datasets
            \begin{itemize}
                \item Each dataset is a separate repository
                \item Output: data files + metadata
            \end{itemize}
        
        \item \textbf{Method modules}: Apply computational methods
            \begin{itemize}
                \item Each method is a separate repository
                \item Input: data from data modules
                \item Output: predictions/results
            \end{itemize}
        
        \item \textbf{Metric modules}: Compute performance measures
            \begin{itemize}
                \item Compare method outputs to ground truth
                \item Generate summary statistics and visualizations
            \end{itemize}
    \end{enumerate}
    
    \vspace{0.3cm}
    \textbf{Central orchestration}: YAML benchmark manifest defines connections and requirements
\end{frame}

\begin{frame}{The clustering example}
    \textbf{Repository:} \url{https://github.com/omnibenchmark/clustering\_example}
    
    \vspace{0.3cm}
    A complete benchmark for clustering methods using the clustbench dataset collection.
    
    \vspace{0.3cm}
    \textbf{Components:}
    \begin{itemize}
        \item \textbf{79+ datasets} from clustbench (mnist, iris, synthetic shapes, etc.)
        \item \textbf{20+ clustering methods} across 5 families:
            \begin{itemize}
                \item Hierarchical (fastcluster, sklearn, genieclust)
                \item Partitioning (k-means, PAM)
                \item Density-based (HDBSCAN)
                \item Model-based (Gaussian mixture)
                \item Specialized (FCPS methods)
            \end{itemize}
        \item \textbf{Multiple metrics}: ARI, MI, FM score, accuracy variants
        \item \textbf{Summary reports}: Comprehensive comparison
    \end{itemize}
    
\end{frame}

\begin{frame}{Clustbench attribution}
    \textbf{Original benchmark by:} Marek Gagolewski
    
    \textbf{Reference:}
    
    Gagolewski M., A framework for benchmarking clustering algorithms, SoftwareX 20, 2022, 101270
    
    \vspace{0.3cm}
    \textbf{Original clustbench:}
    \begin{itemize}
        \item 79 datasets with ground truth labels
        \item Wide variety of cluster shapes and difficulties
        \item Real and synthetic data
        \item Some datasets excluded for computational reasons
    \end{itemize}
    
    \vspace{0.3cm}
    Data repository: \url{https://github.com/imallona/clustbench\_data}
\end{frame}

\begin{frame}[fragile]{Installation}
\begin{lstlisting}[language=bash]
# 1. Install omnibenchmark
pip install omnibenchmark


# 2. Verify installation
ob --version
\end{lstlisting}

\vspace{0.3cm}
\textbf{Full installation guide:}

\url{https://docs.omnibenchmark.org/latest/howto/\#install-omnibenchmark}
\end{frame}

\begin{frame}[fragile]{Cloning the Benchmark}
\begin{lstlisting}[language=bash]
# Clone the clustering example
git clone git@github.com:omnibenchmark/clustering_example.git

# Or using HTTPS
git clone https://github.com/omnibenchmark/clustering_example.git

# Move to the repository
cd clustering_example

# Explore the structure
ls -la
# Key files:
# - Clustering.yaml: Benchmarking manifests
# - envs/: Environment specifications
# - README.md: Documentation
\end{lstlisting}
\end{frame}

\begin{frame}[fragile]{Running the Benchmark}
\begin{lstlisting}[language=bash]
# Basic run (2 cores, local storage)
ob run benchmark -b Clustering.yaml --local-storage --cores 2

# Options:
# -b, --benchmark: Path to configuration YAML
# --local-storage: Store results locally (not S3)
# --cores: Number of parallel processes
# --dry-run: Show what would be executed without running

# Faster (more concurrent tasks)
ob run benchmark -b Clustering.yaml --local-storage --cores 10
\end{lstlisting}

\vspace{0.3cm}
\textit{Warning: Computationally intensive! Start with dry-run.}
\end{frame}

\begin{frame}{Software environments / backends}
    \textbf{Multiple environment backends supported:}
    
    \vspace{0.3cm}
    \begin{itemize}
        \item \textbf{Conda} (default): Python package management
            \begin{itemize}
                \item Config: \texttt{Clustering.yaml}
                \item Best for: Local development, standard packages
            \end{itemize}
        
        \item \textbf{Apptainer/Singularity}: Containers
            \begin{itemize}
                \item Config: \texttt{Clustering\_apptainer.yaml}
                \item Best for: HPC, complex dependencies
            \end{itemize}
        
        \item \textbf{Easybuild/Lmod}: Environment modules
            \begin{itemize}
                \item Config: \texttt{Clustering\_easybuild.yaml}
                \item Best for: HPC with existing modules
            \end{itemize}
    \end{itemize}
    
    \vspace{0.3cm}
    \textbf{See:} \texttt{envs/README.md} for backend details
\end{frame}

\begin{frame}[fragile]{Understanding the Configuration}
\begin{lstlisting}[language=yaml]
# Clustering.yaml (simplified)
benchmark:
  name: "clustering_benchmark"

data:
  - repo: "github.com/imallona/clustbench_data"
    datasets:
      - args: ["--dataset_generator", "mnist", 
               "--dataset_name", "fashion"]
      - args: ["--dataset_generator", "other",
               "--dataset_name", "iris"]

methods:
  - repo: "github.com/imallona/clustbench_fastcluster"
    method:
      - args: ["--linkage", "complete"]
      - args: ["--linkage", "ward"]

metrics:
  - repo: "github.com/imallona/clustbench_metrics"
    metrics:
      - args: ["--metric", "adjusted_rand_score"]
\end{lstlisting}
\end{frame}

\begin{frame}{Data modules}
    \textbf{Data repository:} \url{https://github.com/imallona/clustbench_data}
    
    \vspace{0.3cm}
    \textbf{Dataset generators:}
    \begin{itemize}
        \item \texttt{mnist}: Fashion MNIST, Digits MNIST
        \item \texttt{other}: Iris, chameleon, etc.
        \item \texttt{wut}: Warsaw University of Technology synthetic sets
        \item \texttt{sipu}: Clusters with various shapes
        \item \texttt{uci}: UCI Machine Learning Repository datasets
    \end{itemize}
    
    \vspace{0.3cm}
    \textbf{Each dataset provides:}
    \begin{itemize}
        \item Data matrix (samples × features)
        \item Ground truth labels
        \item Metadata (n\_clusters, difficulty, etc.)
    \end{itemize}
    
    \vspace{0.3cm}
    \textit{Note: Some large datasets commented out for speed}
\end{frame}

\begin{frame}{Method modules}
    \textbf{Five method families:}
    
    \vspace{0.3cm}
    \begin{enumerate}
        \item \textbf{fastcluster} (\url{https://github.com/imallona/clustbench_fastcluster})
            \begin{itemize}
                \item Linkages: complete, ward, average, weighted, median, centroid
            \end{itemize}
        
        \item \textbf{sklearn} (\url{https://github.com/imallona/clustbench_sklearn})
            \begin{itemize}
                \item Methods: birch, kmeans, spectral, Gaussian mixture
            \end{itemize}
        
        \item \textbf{genieclust} (\url{https://github.com/imallona/clustbench_genieclust})
            \begin{itemize}
                \item Methods: genie, GIC, ICA with various thresholds
            \end{itemize}
        
        \item \textbf{FCPS} (\url{https://github.com/imallona/clustbench_fcps})
            \begin{itemize}
                \item Methods: Minimax, MinEnergy, HDBSCAN, Diana, Fanny, etc.
            \end{itemize}
        
        \item \textbf{agglomerative} (\url{https://github.com/imallona/clustbench_agglomerative})
            \begin{itemize}
                \item Linkages: average, complete, ward
            \end{itemize}
    \end{enumerate}
\end{frame}

\begin{frame}{Metric modules}
    \textbf{Metric repository:} \texttt{github.com/imallona/clustbench\_metrics}
    
    \vspace{0.3cm}
    \textbf{External clustering metrics} (require ground truth):
    \begin{itemize}
        \item \texttt{adjusted\_rand\_score}: Adjusted Rand Index (ARI)
        \item \texttt{adjusted\_mutual\_info\_score}: Adjusted MI
        \item \texttt{normalized\_mutual\_info\_score}: Normalized MI
        \item \texttt{fowlkes\_mallows\_score}: FM score
        \item \texttt{normalized\_clustering\_accuracy}: Accuracy variants
        \item \texttt{pair\_sets\_index}: Pair sets index
    \end{itemize}
    
    \vspace{0.3cm}
    \textbf{Each metric provides:}
    \begin{itemize}
        \item Score (0 to 1, higher is better for most)
        \item Runtime and memory usage
        \item Metadata
    \end{itemize}
    
    \vspace{0.3cm}
    \textbf{Summary repository:} \texttt{github.com/imallona/clustering\_report}
\end{frame}

\begin{frame}{Understanding k±2 Strategy}
    \textbf{Challenge:} Different datasets have different true numbers of clusters.
    
    \vspace{0.3cm}
    \textbf{Solution:} Test each dataset with k-2, k-1, k, k+1, k+2 clusters
    
    where k is the true number of clusters.
    
    \vspace{0.3cm}
    \textbf{Special case:} When k=3 (minimum reasonable), use:
    
    k = 2, 2, 3, 5, 6 (fill with k=2 as needed)
    
    \vspace{0.3cm}
    \textbf{Why?}
    \begin{itemize}
        \item Fair comparison: Same number of runs across datasets
        \item Robustness: Tests sensitivity to k parameter
        \item Realism: In practice, true k is often unknown
    \end{itemize}
    
    \vspace{0.3cm}
    \textit{Note: This means some values recomputed, but ensures comparable runtimes}
\end{frame}

\begin{frame}[fragile]{Monitoring Execution}
\begin{lstlisting}[language=bash]
# During execution, monitor progress:
# Check terminal / Snakemake's output

# Check intermediate results
ls -lh out

# Monitor resource usage
top

# Estimate completion time
# Based on jobs completed vs. total

# If execution fails:
# - Check error messages
# - Review log files
# - Verify environment setup
# - Restart with same command (Snakemake resumes)
\end{lstlisting}
\end{frame}

\begin{frame}{Interpreting results}
    \textbf{Key questions to ask:}
    \begin{itemize}
        \item Which methods perform best for most datasets?
        \item Are there dataset-specific strengths/weaknesses?
        \item How much do results vary across datasets?
        \item What's the runtime vs. accuracy tradeoff?
    \end{itemize}
    
    \vspace{0.3cm}
    \textbf{Look for:}
    \begin{itemize}
        \item Methods that consistently rank high
        \item Large performance gaps (potential issues)
        \item Runtime outliers
        \item Failed runs (check logs)
    \end{itemize}
    
\end{frame}

\begin{frame}{Contributing new modules}
    \textbf{Want to add your method, metric or dataset?}
    
    \vspace{0.3cm}
    \textbf{Steps:}
    \begin{enumerate}
        \item Create a new repository for your module
        \item Follow the module contract in terms of inputs, outputs, and arguments
        \item Implement standardized interface:
            \begin{itemize}
                \item Input: data files in expected format
                \item Output: results in expected format
            \end{itemize}
        \item Test locally with toy data
        \item Submit pull request to add module to benchmark YAML
    \end{enumerate}
    
    \vspace{0.3cm}
    \textbf{Documentation:} \texttt{https://docs.omnibenchmark.org/latest/howto/}
    
\end{frame}

\begin{frame}[fragile]{Creating a method module}
\begin{lstlisting}[language=python]
import argparse
import pandas as pd

def main():
    parser = argparse.ArgumentParser()
    parser.add_argument('--data.matrix', required=True)
    parser.add_argument('--data.true_labels', required=True)
    ## etc
    args = parser.parse_args()    

    data = pd.read_csv(args.input)

    labels = cluster(data, n_clusters=args.n_clusters) ## across ks
    
    pd.DataFrame(labels).to_csv(args.output, index=False)

if __name__ == '__main__':
    main()
\end{lstlisting}
\end{frame}

\begin{frame}{Best practices}
    \textbf{Do:}
    \begin{itemize}
        \item Follow standardized interfaces strictly
        \item Pin all dependency versions
        \item Include tests
        \item Document parameters and outputs
        \item Handle errors gracefully
        \item Use appropriate random seeds
    \end{itemize}
    
    \vspace{0.3cm}
    \textbf{Don't:}
    \begin{itemize}
        \item Hardcode paths or parameters
        \item Assume specific working directory
    \end{itemize}
\end{frame}

\begin{frame}{Troubleshooting}
    \textbf{Problem: "Cannot create conda environment"}
    
    Solution: Check internet connection, verify conda installation
    
    \vspace{0.3cm}
    \textbf{Problem: "Method execution failed"}
    
    Solution: Check logs in \texttt{logs/} directory, verify input format
    
    \vspace{0.3cm}
    \textbf{Problem: "Out of memory"}
    
    Solution: Reduce \texttt{--cores}, exclude large datasets, use more RAM
    
    \vspace{0.3cm}
    \textbf{Problem: "Architecture not supported"}
    
    Solution: Use amd64 system (won't work on arm64/M1 Macs)
    
    \vspace{0.3cm}
    \textbf{Problem: "Results differ from expected"}
    
    Solution: Check random seeds, verify dependency versions
\end{frame}


\begin{frame}{Summary}
    \textbf{Key takeaways:}
    \begin{itemize}
        \item Omnibenchmark enables solo and collaborative benchmarking
        \item Modular design: data, methods, metrics as separate repos
        \item Clustering example: Complete benchmark with 79+ datasets, 20+ methods
        \item Running: \texttt{ob run benchmark -b Clustering.yaml --local-storage --cores 6}
        \item Contributing: Create standardized module in new repository
    \end{itemize}
    
    \vspace{0.3cm}
    \textbf{Getting started:}
    \begin{enumerate}
        \item Install omnibenchmark: \texttt{pip install omnibenchmark}
        \item Clone example: \texttt{git clone github.com/omnibenchmark/clustering\_example}
        \item Try dry-run: \texttt{ob run benchmark -b Clustering.yaml --dry-run}
        \item Run
    \end{enumerate}
\end{frame}

\begin{frame}{Resources}
    \begin{itemize}
        \item Omnibenchmark website: \texttt{https://omnibenchmark.org}
        \item Clustering example: \texttt{github.com/omnibenchmark/clustering\_example}
        \item GitHub organization: \texttt{github.com/omnibenchmark}
        \item Contact: Issues on GitHub repositories
    \end{itemize}
    
    \vspace{0.3cm}
    \textbf{Related papers:}
    \begin{itemize}
        \item Gagolewski (2022) "Are cluster validity measures (in)valid?"
    \end{itemize}
\end{frame}

\end{document}
